\documentclass[12pt,a4paper]{article}

\usepackage[utf8]{inputenc}
\usepackage[T1]{fontenc}
\usepackage[french]{babel}
\usepackage[a4paper,margin=2.5cm]{geometry}
\usepackage{amsmath}
\usepackage{amssymb}
\usepackage{booktabs}
\usepackage{array}
\usepackage{hyperref}
\usepackage{xcolor}
\usepackage{listings}

\hypersetup{
  colorlinks=true,
  linkcolor=blue!50!black,
  urlcolor=blue!60!black,
  pdftitle={Fiche de lecture - Risque systémique et Early Warning System},
  pdfauthor={Groupe 10}
}

\lstset{
  basicstyle=\ttfamily\small,
  breaklines=true,
  frame=single,
  columns=fullflexible,
  showstringspaces=false
}

\begin{document}

\begin{titlepage}
  \centering
  \vspace*{1cm}
  {\Large \textbf{Université Cadi Ayyad}\par}
  \vspace{0.3cm}
  {\large École Nationale des Sciences Appliquées de Safi\par}
  \vspace{0.3cm}
  {\large Filière G.INFO et IDIA\par}
  \vfill
  {\Huge \textbf{Fiche de lecture}\par}
  \vspace{0.5cm}
  {\LARGE Risque Systémique \& Early Warning System (CoVaR)\par}
  \vspace{0.8cm}
  {\large Groupe 10\par}
  \vfill
  \begin{tabular}{ll}
    CHAFIK Mouad & ERRAFAI Achraf \\
    MARZAQ Jihane & NADIR Safaa \\
    Encadrant : M. JRAIFI & Année : 2026--2027
  \end{tabular}
  \vfill
\end{titlepage}

\tableofcontents
\clearpage

\section{Objet de la fiche}

Cette fiche présente le cadre mathématique et logiciel du projet
\emph{Risque Systémique \& Early Warning System (CoVaR)}. L'objectif est de mesurer le risque
individuel des banques, leur contribution au risque systémique et les changements atypiques de
la structure de dépendance entre établissements.

Le pipeline actif exploite cinq séries réelles issues des fichiers locaux
\texttt{old-erp/data/raw/} : Attijariwafa Bank (ATW), Banque Centrale Populaire (BCP), Bank of
Africa (BOA), CIH Bank (CIH) et Crédit du Maroc (CDM). Les données sont normalisées dans
\texttt{data/real/}, avec des prix de clôture alignés sur les dates communes. Le jeu synthétique
historique est conservé séparément et n'est pas mélangé aux observations réelles.

\section{Cadre mathématique}

\subsection{Rendement logarithmique}

Pour une banque $i$ et deux dates consécutives, le rendement logarithmique est défini par :
\[
  r_{i,t}=\ln\left(\frac{P_{i,t}}{P_{i,t-1}}\right),
\]
où $P_{i,t}$ désigne le prix de clôture. Les rendements sont calculés après alignement des dates
communes entre les cinq banques.

\subsection{Value at Risk}

La VaR au niveau $q$ est le quantile inférieur de la distribution des rendements :
\[
  \mathbb{P}\left(X_{i,t}\leq \mathrm{VaR}_{q,t}^{i}\right)=q.
\]
Dans l'implémentation historique, la mesure de perte est calculée par
$-\operatorname{quantile}_q(r_i)$. Sous hypothèse normale, la VaR paramétrique s'écrit :
\[
  \mathrm{VaR}_{q}^{i}
  =-\left(\mu_i+\Phi^{-1}(q)\sigma_i\right),
\]
où $\Phi^{-1}$ est le quantile de la loi normale centrée réduite.

\subsection{CoVaR}

La CoVaR mesure le quantile conditionnel du rendement systémique lorsque la banque $i$ est dans
un état donné :
\[
  \mathbb{P}\left(
  X_{\mathrm{sys},t}\leq
  \mathrm{CoVaR}_{q,t}^{\mathrm{sys}|i}
  \,\middle|\,
  X_{i,t}=\mathrm{VaR}_{q,t}^{i}
  \right)=q.
\]
Le rendement systémique est construit comme la moyenne équipondérée des rendements des banques
hors institution étudiée, afin d'éviter une relation mécanique partie--tout.

\subsection{\texorpdfstring{$\Delta$CoVaR}{Delta CoVaR}}

La contribution marginale au risque systémique est :
\[
  \Delta\mathrm{CoVaR}_{t}^{\mathrm{sys}|i}
  =
  \mathrm{CoVaR}_{q,t}^{\mathrm{sys}|i,\text{détresse}}
  -
  \mathrm{CoVaR}_{q,t}^{\mathrm{sys}|i,\text{médiane}}.
\]
Une valeur plus négative indique une détérioration plus forte du rendement systémique lorsque
la banque passe de l'état de référence à l'état de détresse. Cette mesure est une association
conditionnelle et non une preuve de causalité.

\subsection{Régression quantile}

La CoVaR est estimée à partir de la régression quantile :
\[
  X_{\mathrm{sys},t}
  =
  \alpha_q+\gamma_qX_{i,t}+\varepsilon_t.
\]
La prédiction conditionnelle au niveau de la VaR bancaire est alors :
\[
  \mathrm{CoVaR}_{q}^{\mathrm{sys}|i}
  =
  \alpha_q+\gamma_q\mathrm{VaR}_{q}^{i}.
\]
La référence médiane remplace la VaR de détresse par le quantile 50\,\%.

\section{Modèles de volatilité et de dépendance}

\subsection{GARCH(1,1) univarié}

Chaque banque est modélisée séparément par :
\[
  h_{i,t}
  =
  \omega_i+\alpha_i\varepsilon_{i,t-1}^{2}
  +\beta_i h_{i,t-1},
\]
avec :
\[
  \omega_i>0,\qquad
  \alpha_i\geq 0,\qquad
  \beta_i\geq 0,\qquad
  \alpha_i+\beta_i<1.
\]
Les résidus standardisés sont définis par :
\[
  z_{i,t}=\frac{\varepsilon_{i,t}}{\sqrt{h_{i,t}}}.
\]
Les ajustements non convergents ou non admissibles sont rejetés explicitement avant l'étape
multivariée.

\subsection{DCC-GARCH}

Les résidus standardisés des cinq modèles GARCH forment le vecteur $\mathbf{z}_t$. La matrice
de covariance inconditionnelle est :
\[
  \overline{Q}
  =
  \frac{1}{T}\sum_{t=1}^{T}\mathbf{z}_t\mathbf{z}_t^\top.
\]
La récursion DCC est :
\[
  Q_t
  =
  (1-a-b)\overline{Q}
  +a\mathbf{z}_{t-1}\mathbf{z}_{t-1}^{\top}
  +bQ_{t-1},
\]
avec :
\[
  a\geq0,\qquad b\geq0,\qquad a+b<1.
\]
La matrice de corrélation conditionnelle est obtenue par normalisation :
\[
  R_t
  =
  \operatorname{diag}(Q_t)^{-1/2}
  Q_t
  \operatorname{diag}(Q_t)^{-1/2}.
\]
Les matrices $R_t$ sont symétriques, de diagonale unitaire et doivent rester définies positives.
Elles représentent une dépendance conditionnelle dynamique, et non une relation causale.

\section{Détection d'anomalies par autoencodeur}

Pour les cinq banques, la partie triangulaire supérieure stricte de $R_t$ contient :
\[
  \frac{5(5-1)}{2}=10
\]
corrélations uniques. Le vecteur d'entrée est noté $\mathbf{x}_t\in\mathbb{R}^{10}$.

L'encodeur produit une représentation latente :
\[
  \mathbf{z}_t=f_{\theta}(\mathbf{x}_t),
\]
puis le décodeur reconstruit :
\[
  \widehat{\mathbf{x}}_t=g_{\phi}(\mathbf{z}_t).
\]
L'architecture implémentée suit la structure dense :
\[
  10\rightarrow8\rightarrow5\rightarrow3
  \rightarrow5\rightarrow8\rightarrow10,
\]
avec des activations ReLU dans les couches cachées.

L'apprentissage est effectué uniquement sur la fenêtre chronologique considérée comme normale.
Le score d'anomalie quotidien est l'erreur quadratique moyenne :
\[
  s_t
  =
  \frac{1}{10}\sum_{j=1}^{10}
  (x_{t,j}-\widehat{x}_{t,j})^2.
\]
Le seuil d'alerte est le 99e percentile des scores du régime normal :
\[
  \tau_{0.99}=Q_{0.99}(s_t),
  \qquad
  \mathrm{Alerte}_t=\mathbb{1}_{\{s_t>\tau_{0.99}\}}.
\]
Ce signal complète la VaR, la CoVaR et la $\Delta$CoVaR ; il ne remplace pas leur interprétation
économique ou leur validation humaine.

\section{Architecture Odoo et intégration}

Le dépôt contient actuellement l'addon \texttt{cours\_derives/}, issu du TP 0. Les indicateurs
Early Warning issus du pipeline Python pourront être ajoutés dans cet addon ou dans un module
complémentaire dépendant :

\begin{table}[h]
  \centering
  \begin{tabular}{p{0.35\textwidth}p{0.52\textwidth}}
    \toprule
    \textbf{Élément} & \textbf{Rôle} \\
    \midrule
    \texttt{cours.derives.banque} (proposé) &
    Données par banque, ticker, VaR, CoVaR et $\Delta$CoVaR. \\
    \texttt{cours.derives.indicateur} (proposé) &
    Date, corrélation moyenne, CoVaR systémique, score d'anomalie et niveau d'alerte. \\
    \texttt{option\_id} &
    Lien facultatif vers une option du module \texttt{cours\_derives}. \\
    Module \texttt{mail} &
    Support des activités et notifications lors d'une alerte importante. \\
    \bottomrule
  \end{tabular}
  \caption{Structures proposées pour l'intégration Early Warning}
\end{table}

La hiérarchie fonctionnelle peut être organisée ainsi :

\begin{lstlisting}
Derives
|-- Forwards
|-- Options
`-- Early Warning
    |-- Banks
    `-- Indicators
\end{lstlisting}

Le flux de données est :
\[
\text{CSV réels}
\rightarrow
\text{rendements}
\rightarrow
\text{GARCH}
\rightarrow
\text{DCC-GARCH}
\rightarrow
\text{VaR/CoVaR/autoencodeur}
\rightarrow
\text{indicateurs Odoo}.
\]

Chaque alerte doit conserver sa date, son score, le seuil appliqué, le niveau de criticité et la
version du modèle. Cette traçabilité permet de reproduire la décision et de distinguer les
alertes quantitatives des validations métier.

\section{Hypothèses et limites}

\begin{enumerate}
  \item Les cinq banques sont représentées par des rendements de clôture alignés sur les dates
        communes.
  \item Les paramètres GARCH et DCC sont estimés sur l'échantillon disponible ; une évaluation
        prédictive exige une séparation stricte entre calibration et test.
  \item Les corrélations conditionnelles ne démontrent pas une causalité entre établissements.
  \item Le seuil au 99e percentile est empirique et dépend du régime normal choisi.
  \item Une alerte d'autoencodeur signale une reconstruction atypique, mais n'identifie pas
        automatiquement la cause économique du changement.
  \item Les résultats doivent être combinés avec les indicateurs de volatilité, de pertes
        extrêmes et avec l'analyse humaine avant toute décision opérationnelle.
\end{enumerate}

\section{Conclusion}

La chaîne CoVaR combine une modélisation de volatilité individuelle, une dépendance dynamique
DCC-GARCH et une détection non supervisée des anomalies. L'autoencodeur fournit un score continu
sur la structure des corrélations des cinq banques, tandis que le seuil empirique facilite son
intégration dans un tableau de bord Early Warning.

L'étape suivante consiste à automatiser l'import périodique des scores dans
\texttt{cours\_derives.indicateur}, à définir des niveaux vert--orange--rouge et à relier les alertes
aux activités Odoo. Une surveillance de la dérive du modèle devra également être mise en place
pour réestimer le seuil lorsque la structure du marché évoluera durablement.

\end{document}
